\scanpage{135}\mspage{9}
% archive sheet number: 134
La proposition 2 p. 4 de $\unsure{\overline{\underline{\mathrm{XIII}}}|\mathrm{XIII}}$ assure que l'on a
\[
C'''_0\subset\widetilde C_0.
\]
Ceci posé, voici la deuxième propriété essentielle de la factorisation $f=p_1(f)i_1(f)$ (la première étant 1.17.) :
\[
\text{Pour \unsure{toute|tout} }f\in\mathrm{Fl}\,M,\text{ on a}\quad i_1(f)\in C''_0\tag*{(1.20.)}
\]
Il en résulte aussitôt, par récurrence ordinale
\[
\begin{aligned}
&\text{Pour tout }f\in\mathrm{Fl}\,M\text{ et tout ordinal }\alpha\in\unsure{\mathcal U|U},\text{ on a}\\
&i_\alpha(f)\in C'''_0,\quad\text{a fortiori}\quad i_\alpha(f)\in\widetilde C_0.
\end{aligned}\tag*{(1.21)}
\]
Donc si $\alpha$ est grand devant $\unsure{\pi|I}$ (cf cor. 2. p. \deleted{\unsure{0}} 7), la factorisation $f=p_\alpha(f)i_\alpha(f)$ satisfait
\[
\left\{\begin{aligned}
&f=p_\alpha(f)i_\alpha(f),\quad\text{avec }i_\alpha(f)\in\overset{\deleted{\unsure{\sim}}}{C}'''_0\subset\widetilde C_0\inserted{{}=\Phi},\quad p_\alpha(f)\in C_{0*}=\Psi,\\
&\text{lorsque }\alpha\text{ grand devant }\pi\\
&\text{les sources des }i\in C_0\quad\pi\text{-accessibles.}
\end{aligned}\right.\tag*{(1.22.)}
\]
\deleted{\unsure{Pour résumer :}}

\deleted{\uline{Proposition 1} Soit $M$ une $\mathcal U$-catégorie stable par petites $\varinjlim$ filtrantes, $C_0\subset\mathrm{Fl}(M)$ \deleted{\unsure{tel que}}}

Il reste à construire le foncteur $f\mapsto X_1(f)$, avec $i_1(f)$, $p_1(f)$, satisfaisant les conditions (1.17) et (1.20.), sous l'hypothèse (1.16) que les $i\in C_0$ sont \unsure{quarrables}. On considère donc, \inserted{pour $f\in\mathrm{Fl}\,M$ fixé}, l'ens. $\Delta(f)$ \inserted{$(=\Delta(\unsure{C_0|C},f))$} de \unsure{tous} les diagrammes $D$ de la forme 1.15., \inserted{avec $i\in C_0$}, et on pose
\[
X_1(f)=\coprod_{D\in\Delta(f)\,(X)}X_D\tag*{(1.23)}
\]
(somme amalgamée \unsure{sous|dans} $X$, i.e. somme \unsure{indiciée|indicée|inductive|induite}
